% constraints-risks-checklist.tex
% Document 5 - Constraints, Risks, and Compliance Checklist
%
% Build (from this folder, requires XeLaTeX -- see README.md):
%   xelatex -interaction=nonstopmode constraints-risks-checklist.tex
%   xelatex -interaction=nonstopmode constraints-risks-checklist.tex   (rerun for TOC/page refs)
%
% Completed with the requesting operations/safety/IT representatives and
% the center lead, typically alongside the Scope Definition Worksheet
% (Document 4). See docs/Engineering_Project_Intake_Plan.md for the full
% workflow.

\documentclass[11pt]{article}

\usepackage{../../shared/intake-style}
\usepackage{../../shared/form-components}

\SetDocIdentity{Constraints, Risks, and Compliance Checklist}{ENG-FORM-05}{v0.1-draft}
\SetFieldPrefix{constrisk}

% Replace these placeholders once branding and markings are confirmed
% (see "Remaining Organization-Specific Decisions" in the intake plan).
\SetCompanyName{[Engineering Center Name]}
\SetMarkingText{FOR OFFICIAL USE ONLY}

\SetDocPurpose{%
This checklist records the schedule constraints, applicable codes and
standards, safety and environmental constraints, data-handling
restrictions, and known risks that could affect this work. Its purpose is
to surface these things early, while they are still easy to plan around,
rather than discovering them after work has started. Recording an
applicable standard here means it has been identified as relevant -- it
does not mean the work has been certified as compliant with it. That
determination is made separately by the appropriate discipline lead or
compliance authority.%
}

\begin{document}

\begin{Form}

\CoverPage

\DocumentControlBlock

\LinkedIntakeBlock

\RequestingOrgBlock

\IntakeSection{Schedule Constraints}
\TextArea{SchedConstraints}{Known outage windows, access windows, blackout periods, or other scheduling constraints}{\linewidth}{1.5cm}

\IntakeSection{Applicable Codes and Standards}
\Guidance{List applicable codes/standards across all disciplines involved.
This records applicability and who confirmed it, not a compliance
certification.}\\[4pt]
\StandardsTable{ComplStd}

\IntakeSection{Safety Constraints}
\TextArea{SafetyConstraints}{Known safety constraints, required PPE, permits-to-work, or safety reviews}{\linewidth}{1.5cm}

\IntakeSection{Environmental Constraints}
\TextArea{EnvConstraints}{Known environmental constraints, permits, or site-condition limitations}{\linewidth}{1.5cm}

\IntakeSection{Data and Information-Handling Restrictions}
\TextArea{DataRestrictions}{Known data classification, export-control, or information-handling restrictions that apply to this work}{\linewidth}{1.5cm}

\IntakeSection{Risk Register}
\Guidance{List known risks to the schedule, scope, or safety of this
work, however early-stage.}\\[4pt]
\RiskRegisterTable{Risk}

\IntakeSection{Specialist Reviews Needed}
\Guidance{List any reviews by a specialist, safety officer, or other
authority that this work is known to require.}\\[4pt]
\begin{tabularx}{\textwidth}{Y{0.3}Y{0.3}Y{0.2}Y{0.2}}
  \toprule
  Review Type & Reviewer / Organization & Status & Date \\
  \midrule
  \TextLine{Rev1Type}{}{\linewidth} & \TextLine{Rev1Org}{}{\linewidth} & \TextLine{Rev1Status}{}{\linewidth} & \TextLine{Rev1Date}{}{\linewidth} \\[8pt]
  \TextLine{Rev2Type}{}{\linewidth} & \TextLine{Rev2Org}{}{\linewidth} & \TextLine{Rev2Status}{}{\linewidth} & \TextLine{Rev2Date}{}{\linewidth} \\[8pt]
  \TextLine{Rev3Type}{}{\linewidth} & \TextLine{Rev3Org}{}{\linewidth} & \TextLine{Rev3Status}{}{\linewidth} & \TextLine{Rev3Date}{}{\linewidth} \\
  \bottomrule
\end{tabularx}\\[4pt]
\ContinuationNote

\IntakeSection{Assumptions and Open Questions}
\OpenItemsTable{ComplOpen}

\AcknowledgementBlock{This checklist records known constraints, applicable
standards, and risks at the time of signing. It does not certify
compliance with any listed standard, does not commit \CompanyName's
resources, and is not authorization to begin engineering work or to
charge project time.}

\end{Form}

\end{document}
